\begin{proof}[Proof of \cref{obj:lem:chebyshev-factorial-certificate}]
\begin{enumerate}
\item The Chebyshev recurrence gives
\[
T_h(1)=1,\qquad \deg T_h\le h,\qquad
T_h(\cos\theta)=\cos(h\theta)
\quad(h\ge0,\ 0\le\theta\le\pi).
\]
Indeed, the initial values and the cosine addition formula establish the last identity by induction. Consequently,
\[
|T_L(x)|\le1\qquad(-1\le x\le1).
\]
Applying \cref{obj:lem:markov-polynomial} with polynomial \(T_L\), degree bound \(L\), and uniform bound \(1\), we obtain
\[
|T_L'(x)|\le L^2\qquad(-1\le x\le1).
\]
For \(0\le z\le1\), the points \(1-2z\) and \(1\) both belong to this interval. The derivative bound therefore gives
\[
|T_L(1-2z)-T_L(1)|\le L^2|(1-2z)-1|=2L^2z.
\]
Combining this with the range bound and \(T_L(1)=1\) yields
\[
0\le1-C_L(z),\qquad
1-C_L(z)\le2,\qquad
1-C_L(z)\le2L^2z.
\]

The numerator \(1-C_L\) has zero constant coefficient. Polynomial division by \(z\) thus gives the identity
\[
zE_L(z)=\frac{1-C_L(z)}{2L^2}.
\]
For \(z\ne0\), this implies
\[
E_L(z)=\frac{1-C_L(z)}{2L^2z}.
\]
When \(0<z\le1\), the denominator is positive. Dividing the three numerator bounds by it gives
\[
0\le E_L(z),\qquad
E_L(z)\le1,\qquad
E_L(z)\le\frac{1}{L^2z},
\]
and hence the required minimum bound.
% lean: CausalSmith.Stat.AnnotationRarearmFrontier.chebE_approximation_bounds

\item Differentiating the polynomial identity from the preceding step and evaluating at zero gives
\[
E_L(0)=-\frac{C_L'(0)}{2L^2}
      =\frac{T_L'(1)}{L^2}.
\]
The endpoint derivative identity is
\[
T_L'(1)=L^2.
\]
It follows directly by differentiating the Chebyshev recurrence: the initial derivatives are
\[
T_0'(1)=0,\qquad T_1'(1)=1,
\]
and, using \(T_h(1)=1\),
\[
T_{h+2}'(1)=2+2T_{h+1}'(1)-T_h'(1).
\]
Induction gives \(T_h'(1)=h^2\). Since \(L\ge2\), division by \(L^2\) is valid, and therefore
\[
E_L(0)=1.
\]
Thus \(1-E_L\) also has zero constant coefficient. Its polynomial division by \(z\) gives
\[
zG_L(z)=1-E_L(z).
\]
For every \(z\ne0\), the two polynomial identities consequently yield
\[
E_L(z)=\frac{1-C_L(z)}{2L^2z},
\qquad
G_L(z)=\frac{1-E_L(z)}{z}.
\]
These identities specify the polynomial continuations at zero as well.
% lean: CausalSmith.Stat.AnnotationRarearmFrontier.cheb_continuation_eval

\item Composition with the affine polynomial \(1-2z\) gives
\[
\deg C_L\le L,\qquad \deg(1-C_L)\le L.
\]
Scalar multiplication preserves this degree bound, and polynomial division by \(z\) lowers it by one. Applying the same argument to \(1-E_L\) gives
\[
\deg E_L\le L-1,\qquad
\deg(1-E_L)\le L-1,\qquad
\deg G_L\le L-2.
\]
The inequalities also hold for a zero polynomial, with the usual degree convention. Therefore the coefficient expansion is
\[
G_L(z)=\sum_{h=0}^{L-2}g_hz^h.
\]
% lean: CausalSmith.Stat.AnnotationRarearmFrontier.chebG_natDegree_le

\item For any real polynomial \(p_{\mathrm{aux}}\), use coefficient notation and define its coefficient norm by
\[
p_{\mathrm{aux}}(z)
 =\sum_{h\ge0}\bigl([z^h]p_{\mathrm{aux}}\bigr)z^h,
\qquad
\|p_{\mathrm{aux}}\|_{\mathrm{coef},1}
 :=\sum_{h\ge0}|[z^h]p_{\mathrm{aux}}|.
\]
Both sums have finite support. The coefficient triangle inequality and convolution formula give, for real polynomials \(p_{\mathrm{aux}},q_{\mathrm{aux}}\) and \(a_{\mathrm{sc}}\in\mathbb R\),
\[
\begin{aligned}
\|p_{\mathrm{aux}}-q_{\mathrm{aux}}\|_{\mathrm{coef},1}
 &\le \|p_{\mathrm{aux}}\|_{\mathrm{coef},1}
      +\|q_{\mathrm{aux}}\|_{\mathrm{coef},1},\\
\|p_{\mathrm{aux}}q_{\mathrm{aux}}\|_{\mathrm{coef},1}
 &\le \|p_{\mathrm{aux}}\|_{\mathrm{coef},1}
       \|q_{\mathrm{aux}}\|_{\mathrm{coef},1},\\
\|a_{\mathrm{sc}}p_{\mathrm{aux}}\|_{\mathrm{coef},1}
 &=|a_{\mathrm{sc}}|\|p_{\mathrm{aux}}\|_{\mathrm{coef},1}.
\end{aligned}
\]
Multiplication by \(z\) preserves this norm:
\[
\|zp_{\mathrm{aux}}\|_{\mathrm{coef},1}
 =\|p_{\mathrm{aux}}\|_{\mathrm{coef},1}.
\]
For the zero polynomial both sides are zero; otherwise this follows by shifting the coefficients, since
\[
[z^0](zp_{\mathrm{aux}})=0,\qquad
[z^{h+1}](zp_{\mathrm{aux}})=[z^h]p_{\mathrm{aux}}.
\]

Define the shifted family
\[
C_{\mathrm{sh},h}(z):=T_h(1-2z)
\quad(h\ge0,\ z\in\mathbb R),
\qquad C_{\mathrm{sh},L}=C_L.
\]
Its initial coefficient norms and recurrence are
\[
\|C_{\mathrm{sh},0}\|_{\mathrm{coef},1}=1,\qquad
\|C_{\mathrm{sh},1}\|_{\mathrm{coef},1}=3,
\]
\[
C_{\mathrm{sh},h+2}
 =2(1-2z)C_{\mathrm{sh},h+1}-C_{\mathrm{sh},h}.
\]
Since
\[
\|2(1-2z)\|_{\mathrm{coef},1}=|2|+|-4|=6,
\]
the norm inequalities imply
\[
\|C_{\mathrm{sh},h+2}\|_{\mathrm{coef},1}
 \le6\|C_{\mathrm{sh},h+1}\|_{\mathrm{coef},1}
    +\|C_{\mathrm{sh},h}\|_{\mathrm{coef},1}.
\]
Starting from \(1\le7^0\) and \(3\le7^1\), two-step induction gives
\[
\|C_{\mathrm{sh},h}\|_{\mathrm{coef},1}\le7^h,
\]
because
\[
6\cdot7^{h+1}+7^h
 =43\cdot7^h
 \le49\cdot7^h
 =7^{h+2}.
\]

Applying the norm identities to \(zE_L=(1-C_L)/(2L^2)\), we obtain
\[
\|E_L\|_{\mathrm{coef},1}
 =\frac{\|1-C_L\|_{\mathrm{coef},1}}{2L^2}
 \le\frac{1+7^L}{2L^2}
 \le7^L.
\]
The last inequality uses \(L^2\ge1\) and \(7^L\ge1\). Applying them next to \(zG_L=1-E_L\) yields
\[
\begin{aligned}
\sum_{h=0}^{L-2}|g_h|
 &=\|G_L\|_{\mathrm{coef},1}\\
 &=\|1-E_L\|_{\mathrm{coef},1}\\
 &\le1+\|E_L\|_{\mathrm{coef},1}\\
 &\le1+7^L
 \le2\cdot7^L
 \le2^L7^L
 =14^L.
\end{aligned}
\]
Here the first equality uses the degree bound, and \(2\le2^L\) follows from \(L\ge2\).
% lean: CausalSmith.Stat.AnnotationRarearmFrontier.chebG_coeff_sum_le

\item Fix \(t,B>0\), \(z\ge0\), and \(D=tB\ge L\), and define the Poisson mean
\[
\lambda_{\mathrm P}:=Dz=tzB\ge0.
\]
The law of \(K\) is
\[
\mathbb P(K=k)=e^{-\lambda_{\mathrm P}}
              \frac{\lambda_{\mathrm P}^{\,k}}{k!}
\qquad(k\in\mathbb N_0).
\]
We use \(0^0=1\) and the empty-product value \(1\).

For every integer \(h\ge0\), the terms with \(k<h\) have \((k)_h=0\), while \((k)_h=k!/(k-h)!\) for \(k\ge h\). Shifting the Poisson series therefore gives
\[
\begin{aligned}
\mathbb E[(K)_h]
 &=e^{-\lambda_{\mathrm P}}
   \sum_{k=h}^{\infty}
   \frac{\lambda_{\mathrm P}^{\,k}}{(k-h)!}\\
 &=\lambda_{\mathrm P}^{\,h}e^{-\lambda_{\mathrm P}}
   \sum_{k=0}^{\infty}\frac{\lambda_{\mathrm P}^{\,k}}{k!}
 =\lambda_{\mathrm P}^{\,h}.
\end{aligned}
\]
The exponential series converges, so these factorial moments are finite. This computation includes \(\lambda_{\mathrm P}=0\) and \(h=0\).

Since \(D>0\), each summand of the estimator satisfies
\[
\mathbb E\!\left[g_h\frac{(K)_h}{D^h}\right]
 =g_h\frac{(Dz)^h}{D^h}
 =g_hz^h.
\]
Integrating the finite sum term by term and using the polynomial expansion gives
\[
\mathbb E[\widehat G]
 =\sum_{h=0}^{L-2}g_hz^h
 =G_L(z).
\]
% lean: CausalSmith.Stat.AnnotationRarearmFrontier.chebG_factorial_unbiased

\item To obtain the mixed moments, define the falling-factorial polynomials
\[
F_{\mathrm{fac},h}(x):=\prod_{i=0}^{h-1}(x-i)
\quad(h\in\mathbb N_0,\ x\in\mathbb R),
\qquad F_{\mathrm{fac},0}(x)=1.
\]
For an integer \(k\ge0\), they satisfy
\[
F_{\mathrm{fac},h}(k)=(k)_h.
\]

Let \(h,h'\in\mathbb N_0\). First suppose \(h\le h'\). The binomial addition identity for falling-factorial polynomials, applied to \(x=(x-h)+h\), gives
\[
F_{\mathrm{fac},h'}(x)
 =\sum_{j_{\mathrm{ov}}=0}^{h'}
   \binom{h'}{j_{\mathrm{ov}}}
   F_{\mathrm{fac},j_{\mathrm{ov}}}(h)
   F_{\mathrm{fac},h'-j_{\mathrm{ov}}}(x-h).
\]
The product definition supplies
\[
F_{\mathrm{fac},h}(x)
F_{\mathrm{fac},h'-j_{\mathrm{ov}}}(x-h)
 =F_{\mathrm{fac},h+h'-j_{\mathrm{ov}}}(x).
\]
Moreover,
\[
F_{\mathrm{fac},j_{\mathrm{ov}}}(h)=0
 \quad(j_{\mathrm{ov}}>h),
\qquad
F_{\mathrm{fac},j_{\mathrm{ov}}}(h)
 =j_{\mathrm{ov}}!\binom{h}{j_{\mathrm{ov}}}
 \quad(0\le j_{\mathrm{ov}}\le h).
\]
Multiplying the addition identity by \(F_{\mathrm{fac},h}(x)\) and removing its zero terms thus gives
\[
F_{\mathrm{fac},h}(x)F_{\mathrm{fac},h'}(x)
 =\sum_{j_{\mathrm{ov}}=0}^{h}
   \binom{h}{j_{\mathrm{ov}}}
   \binom{h'}{j_{\mathrm{ov}}}
   j_{\mathrm{ov}}!\,
   F_{\mathrm{fac},h+h'-j_{\mathrm{ov}}}(x).
\]
If \(h'>h\) fails, exchange the two orders; the product and the coefficients are symmetric. Evaluating the resulting identity at any integer \(k\ge0\) yields
\[
(k)_h(k)_{h'}
 =\sum_{j_{\mathrm{ov}}=0}^{\min(h,h')}
   \binom{h}{j_{\mathrm{ov}}}
   \binom{h'}{j_{\mathrm{ov}}}
   j_{\mathrm{ov}}!\,(k)_{h+h'-j_{\mathrm{ov}}}.
\]
% lean: Causalean.Stat.Concentration.Poisson.descFactorial_mul

\item Define, for \(h,h'\in\mathbb N_0\),
\[
M_{\mathrm{fac}}(h,h')
 :=\mathbb E\!\left[
       \frac{(K)_h}{D^h}\frac{(K)_{h'}}{D^{h'}}
     \right].
\]
The product expansion and the finite factorial moments justify termwise expectation. They give
\[
\begin{aligned}
M_{\mathrm{fac}}(h,h')
 &=\sum_{j_{\mathrm{ov}}=0}^{\min(h,h')}
   \binom{h}{j_{\mathrm{ov}}}
   \binom{h'}{j_{\mathrm{ov}}}
   j_{\mathrm{ov}}!\,
   \frac{(Dz)^{h+h'-j_{\mathrm{ov}}}}{D^{h+h'}}\\
 &=\sum_{j_{\mathrm{ov}}=0}^{\min(h,h')}
   \binom{h}{j_{\mathrm{ov}}}
   \left(
     \frac{\binom{h'}{j_{\mathrm{ov}}}j_{\mathrm{ov}}!}
          {D^{j_{\mathrm{ov}}}}
   \right)
   z^{h+h'-j_{\mathrm{ov}}}.
\end{aligned}
\]
All terms are nonnegative, so \(M_{\mathrm{fac}}(h,h')\ge0\).

Now suppose \(h,h'\le L\). Since \(h'\le L\le D\), the overlap coefficient satisfies
\[
0\le
\frac{\binom{h'}{j_{\mathrm{ov}}}j_{\mathrm{ov}}!}
     {D^{j_{\mathrm{ov}}}}
 =\frac{(h')_{j_{\mathrm{ov}}}}{D^{j_{\mathrm{ov}}}}
 \le\frac{(h')^{j_{\mathrm{ov}}}}{D^{j_{\mathrm{ov}}}}
 \le1.
\]
The falling-factorial bound follows by bounding each factor by \(h'\); when the order exceeds \(h'\), the falling factorial is zero. The order-zero case has value \(1\).

Also, for every overlap index in the displayed sum,
\[
z^{h+h'-j_{\mathrm{ov}}}
 \le(1+z)^{h+h'-j_{\mathrm{ov}}}
 \le(1+z)^{2L},
\]
because \(z\ge0\), \(1+z\ge1\), and \(h+h'-j_{\mathrm{ov}}\le2L\). These inequalities include \(z=0\) under the stated power convention. Hence
\[
\begin{aligned}
M_{\mathrm{fac}}(h,h')
 &\le\sum_{j_{\mathrm{ov}}=0}^{\min(h,h')}
       \binom{h}{j_{\mathrm{ov}}}(1+z)^{2L}\\
 &\le\sum_{j_{\mathrm{ov}}=0}^{h}
       \binom{h}{j_{\mathrm{ov}}}(1+z)^{2L}\\
 &=2^h(1+z)^{2L}\\
 &\le2^L(1+z)^{2L}.
\end{aligned}
\]
The second inequality adds nonnegative terms, and the equality is the binomial theorem.
% lean: CausalSmith.Stat.AnnotationRarearmFrontier.normalized_factorial_mixed_bound

\item Expanding the square and integrating the finite sum gives
\[
\mathbb E[\widehat G^2]
 =\sum_{h=0}^{L-2}\sum_{h'=0}^{L-2}
   g_hg_{h'}M_{\mathrm{fac}}(h,h').
\]
Every product is integrable by its finite factorial expansion. For each pair of indices,
\[
g_hg_{h'}\le|g_h||g_{h'}|,
\qquad
0\le M_{\mathrm{fac}}(h,h')
 \le2^L(1+z)^{2L},
\]
where the moment bound applies because \(h,h'\le L-2\le L\). Consequently,
\[
\begin{aligned}
\mathbb E[\widehat G^2]
 &\le\sum_{h=0}^{L-2}\sum_{h'=0}^{L-2}
       |g_h||g_{h'}|\,2^L(1+z)^{2L}\\
 &=\left(\sum_{h=0}^{L-2}|g_h|\right)^2
      2^L(1+z)^{2L}\\
 &\le(14^L)^2\,2^L(1+z)^{2L}\\
 &=392^L(1+z)^{2L}.
\end{aligned}
\]
The coefficient sum is nonnegative, so its bound may be squared. All index ranges remain valid at \(L=2\), when the estimator sum consists of its order-zero term.
% lean: CausalSmith.Stat.AnnotationRarearmFrontier.chebG_factorial_second_moment
\end{enumerate}
\end{proof}
